% frontmatter/introduccion.tex
\chapter*{Introducci\'on}
\addcontentsline{toc}{chapter}{Introducci\'on}
\label{chapter:introduccion}

La transformaci\'on digital ha dejado de ser una ventaja diferencial para convertirse en una condici\'on
operativa de las organizaciones contempor\'aneas. En las organizaciones sin fines de lucro y de
voluntariado (ONG), la gesti\'on de personas, proyectos, recursos y procesos de gobernanza suele
apoyarse en herramientas fragmentadas —hojas de c\'alculo, mensajer\'ia informal y formularios
dispersos— que dificultan la trazabilidad, comprometen la seguridad de datos sensibles y limitan la
rendici\'on de cuentas ante donantes y entes reguladores. El modelo de software como servicio (SaaS)
multi-tenant responde a esta necesidad al ofrecer, desde una \'unica base de c\'odigo, un servicio
compartido por m\'ultiples organizaciones, manteniendo el aislamiento estricto de sus datos
\citep{fowler2015, redhat2023}.

El presente informe describe y eval\'ua \textbf{Democra}, una plataforma web SaaS multi-tenant para la
gesti\'on integral de ONG, tal como se encuentra materializada en su repositorio de c\'odigo fuente. El
sistema —originalmente desarrollado como \textit{SistemaVolV2.0} e integrado posteriormente en el
\textit{monorepo} Democra— implementa una arquitectura desacoplada de cliente y servidor, un modelo
de datos relacional multi-esquema sobre PostgreSQL administrado a trav\'es de Supabase, y un conjunto
de mecanismos transversales de seguridad, auditor\'ia y control de acceso, gobernados por la
metodolog\'ia de documentaci\'on de dise\~no DDS.

La documentaci\'on se elabora bajo un enfoque de \textit{ingenier\'ia inversa documental}: el repositorio
constituye la \'unica fuente de verdad, y el an\'alisis describe fielmente lo que el c\'odigo implementa,
sin inventar funcionalidades, arquitecturas ni m\'etricas inexistentes. Las secciones cuyos valores
dependen de la ejecuci\'on de pruebas en un entorno controlado —cobertura efectiva, tiempos de
respuesta, resultados de usabilidad— se marcan expl\'icitamente como pendientes de evidencia.

El informe adopta la estructura IMRyD (Introducci\'on, M\'etodo, Resultados y Discusi\'on) propia de la
investigaci\'on aplicada. El primer cap\'itulo plantea el problema, la hip\'otesis, los
objetivos y la justificaci\'on. El \autoref{chapter:marco} desarrolla los antecedentes, el marco
te\'orico-tecnol\'ogico, la metodolog\'ia DDS y la justificaci\'on del \textit{stack}. El
\autoref{chapter:metodo} expone el material y los m\'etodos. El \autoref{chapter:arquitectura} describe
la arquitectura. El \autoref{chapter:datos} detalla el modelo de datos. El \autoref{chapter:modulos}
documenta los m\'odulos funcionales. El \autoref{chapter:api} especifica la API y los servicios de
seguridad. El \autoref{chapter:calidad} expone el aseguramiento de la calidad. El
\autoref{chapter:resultados} presenta los resultados y el \autoref{chapter:discusion} su discusi\'on.
El \autoref{chapter:conclusiones} sintetiza conclusiones y recomendaciones.
